% ============================================================
% Abstract (English)
% ============================================================
\chapter*{Abstract}

The contemporary practice of API quality assurance is fragmented across loosely connected concerns: structural validation linters check whether an OpenAPI document parses, contract-testing tools verify that the API server matches its specification at a few hand-chosen endpoints, and load-testing harnesses measure performance under hand-authored scenarios. Each tool is independently useful; none addresses the underlying problem, which is that the specification itself is the only artefact common to development, testing, deployment, and operations, and that its quality determines the quality of every downstream automation that depends on it. This thesis proposes a different approach.

The proposal is developed as a three-layer contribution package: \textbf{SDT}, a conceptual framework that re-grounds API quality assurance in the specification itself; \textbf{DriveBy}, the reference implementation that makes the framework executable; and \textbf{XSDLC}, the GitOps integration that embeds it into the delivery pipeline as a declaratively provisioned quality gate. Each layer is summarised in turn below, followed by the empirical evaluation.

\textbf{Specification-Driven Testing (SDT)} is a conceptual framework for automated API quality assurance in which the specification is treated as the canonical source of truth for both \emph{validation} (does the contract describe a well-formed, complete, secure API?) and \emph{testing} (does the implementation satisfy the contract under realistic functional and load conditions?). The framework is anchored on three axioms---\emph{Completeness}, \emph{Determinism}, and \emph{Observability}---and decomposed into nine principles (P001 through P009) covering structural compliance, documentation quality, error-handling, schema precision, security, functional correctness, performance, versioning, and test-readiness. The contribution at this layer is theoretical: a framework, not a complete methodology; the missing methodology elements (formal test-derivation procedure, execution and failure semantics, lifecycle governance) are scoped as future work.

\textbf{DriveBy} is the reference implementation of the framework: a Go command-line validator that accepts OpenAPI 3.0.x and 3.1.0 specifications natively---and Swagger 2.0 documents through a transparent conversion adapter---runs the seven currently implemented principle checkers (P001--P005, P008, P009; P006 and P007 are scoped to the runtime test layer), and emits a deterministic, machine-readable JSON report with severity-classified findings. The CLI is built around an \texttt{APISpec} adapter that abstracts away version-specific details, allowing every validator to operate on a single semantic surface regardless of whether the underlying spec is OpenAPI 3.0.x, 3.1.0, or a converted Swagger 2.0 document.

\textbf{XSDLC} (the Crossplane Software Development Lifecycle composition) is the GitOps integration: a single Kubernetes custom resource that declaratively provisions a complete two-repository GitOps delivery pipeline---application repository, dedicated GitOps repository, ArgoCD applications with source-hydration overlays, the GitOps Promoter, multi-stage quality gates running DriveBy at staging and production boundaries, branch protection, and the ingress and certificate scaffolding required to run the system. XSDLC is the empirical artefact that demonstrates SDT can be embedded as a built-in feature of an Internal Developer Platform rather than a bolt-on linter.

The framework, the CLI, and the composition are evaluated empirically across four arms. A \emph{controlled evaluation} on a hand-crafted reference API (\texttt{non-critical-api}) verifies defect detection at the level of individual checks: 12 of 15 injected single-check mutations are detected directly (an 80\% detection rate; of the remainder, one mutation is neutralised by the loader's structural validation, one proves to be an injection artefact with no real defect, and exactly one is a genuine framework miss, documented as a limitation), and all 3 evaluable random multi-principle combinations are fully detected, every targeted principle flagged simultaneously. A \emph{large-scale evaluation} on 50 OpenAPI 3.x specifications harvested from the APIs.guru registry shows that 84\% of public APIs score $1/6$ or worse in DriveBy strict mode---no API in the sample passes more than 2 of the 6 evaluated principles. An \emph{operational proof-of-concept} deploys XSDLC on Novelcore's internal Kubernetes cluster and processes five APIs through the full delivery pipeline. Finally, an \emph{agent-feedback experiment} on the canonical Swagger Petstore specification demonstrates that an AI coding agent, given only the DriveBy report as feedback, can autonomously revise a 1/6-scoring specification to 4/6 in a single iteration---empirical evidence for the framework's positioning as feedback infrastructure for agent-driven specification improvement.

The thesis closes with an explicit treatment of authorship and the development methodology, separating the author's intellectual ownership of the framework, architecture, and evaluation design from the AI-assisted productisation phase, with quantitative evidence drawn from the version-control history.

\noindent\textbf{Keywords:} API quality assurance, specification-driven testing, OpenAPI, Crossplane, GitOps, Kubernetes, quality gates, agent-driven development.

\clearpage
